\clearpage
\section*{September 9, 2026: The full parent review finds duplicate proposals and missing buildings}
\textit{Agent-drafted audit; recommendations have not changed production membership.}

The review covers all 76 multi-constituent parents in the current estimation sample: 25 historical and 51 post-policy parents, containing 168 filing constituents. Each has a written assessment, with varying corroboration; not every connection or unit count is independently verified.

Five parents add withdrawn applications to their replacements. Six old applications were withdrawn before new applications appeared with matching building identification numbers, owners and architects. Removing the old applications reduces the recorded totals as follows:

\begin{center}
\begin{tabular}{lrr}
\hline
Site & Recorded units & Replacement proposal \\
\hline
Hillside Avenue / Parsons Boulevard & 350 & 154 \\
957 East 232nd Street & 70 & 35 \\
2279 University Avenue & 96 & 48 \\
1609 Crosby Avenue & 72 & 36 \\
176 East 206th Street & 106 & 53 \\
\hline
\end{tabular}
\end{center}

These six applications add 368 units and create four apparent multi-building parents. Three would fall below the 50-unit estimation cutoff when treated as single buildings. The appropriate cohort date remains a separate choice. Three additional links need stronger evidence: 767/769 East 232nd Street, 1660/1674 Boone Avenue, and 3367 Wilson Avenue/3374 Boston Road have different named sponsors and architects. Parcel adjacency or old lot history alone does not establish common development.

The nearby screen checks filings within 200 metres and 365 days of any constituent, plus same-lot matches regardless of distance. Across all 876 estimation parents, including singletons, it finds 1,238 parent--filing pairs. The 137 pairs around the 76 multi-constituent parents received an initial record review. Seven merit follow-up as missing companions: 290 Bergen Street, 278 Butler Street, 139 Beach 29th Street, 35-11 42nd Street, 4137 Third Avenue, 8 Stuyvesant Place and 3 West Street. Butler explicitly lists all three development lots. The remaining 1,101 comparisons around singletons are unadjudicated. Five anchor filings lack coordinates; the screen cannot establish completeness outside its bounds.

\paragraph{Update to the Third Avenue finding.} Later foundation and structural DOB records also report 99 units at both 4121 and 4133 Third Avenue; the Housing Database's 99+99 has additional administrative support. The saved I1 records still say 98+97. Moreover, 4137 Third Avenue has a nearby 99-unit filing with the same owner and applicant, filed the same day as 4133. A three-building boundary would imply 297 units using the Housing Database-priority measure or 294 using saved I1 counts. Neither is an approved decision-date unit schedule. A separate sale of 4121 in May 2026 also makes the ownership date consequential.

Before estimation, resolve replacements, uncertain boundaries and unit definitions. Common sponsorship does not establish separate wage assessment.

\textit{Reproduction note.} Built with \texttt{make} in
\path{tasks/audits/audit_estimation_parent_links/code/} against revision
\texttt{8458668}, HDB 25Q4 and the saved July 2026 DOB capture. See \path{output/parent_casebook.md} for all cases, sources and nearby notes.
The data reports establish the first audit baseline. No production panel or estimate changed.
